\scanpage{109}
\mspage{105}
% archive sheet number: 108

\uline{Lemme 1} Soit $b\leq c$ dans $\unsure{\Sigma|\bar\Sigma}$. Conditions équivalentes
\begin{enumerate}[label={},leftmargin=1.8em]
\item (i) $\exists b^*\in\Sigma$, avec $b^*\,|\unsure{\circ|0|\bullet}|\,b$ et $b\vee b^*=c$
\item (ii) On a $b\vee b_0^*=c$, où $b_0^*=\unsure{C|\complement}b\wedge c$
\end{enumerate}
De plus, si on a $\unsure{b'|b^*}$ comme en (i), on a $b^*\leq b_0^*$, i.e. $b_0^*$ est le plus grand des $b^*$ satisfaisant (i).

\uline{Lemme 2} Soit $b\leq c$ dans $\Sigma$. Conditions équivalentes
\begin{enumerate}[label={},leftmargin=1.8em]
\item (i) $b,c$ engendrent dans $\Sigma$ \inserted{(pour $\vee,\wedge,C$)} une sous-algèbre de Boole.
\item (ii) $b,c$ engendrent une sous-algèbre de Boole finie, de cardinal $2^n$ avec $n\in\{1,2,3\}$.

\inserted{NB} (($n\deleted{\unsure{\leq}}=1$ sss $\{b,c\}\subset\{0,1\}$ ; $n\leq2$ sss $b=0$ ou $c=1$ ou $b=c$ ; \deleted{mais $\{b,c\}\ne\{0,1\}$} \inserted{\deleted{$\{b,c\}\ne\{0,1\}$ (donc exclus}} \deleted{\unsure{mais} $c\ne0$, $b\ne1$ et les cas $b=\unsure{c|b}=0$, $b=\unsure{c|b}=1$, \unsure{ou} \inserted{et} $b=0$ et $c=1$)} ; $n=3$ sss $0<b<c<1$).
\item \unsure{(iii)|(iii')} $(b,c)$ satisfont les conditions du lemme 1, et $(c',b')$ aussi (où $c'=Cc$, $b'=Cb$).
\item \inserted{(iv) Posant $\alpha=b$, $\beta=c\wedge Cb$, $\unsure{\gamma|\sigma}=Cc$, \unsure{les|b,} $\alpha,\beta,\gamma$ engendrent une sous-alg. de Boole \deleted{\illegible{1}} \unsure{ou encore, on a les} \deleted{\unsure{aux}} relations $C(\alpha\vee\beta)=\gamma$, $C(\beta\vee\gamma)=\alpha$ \inserted{$C(\gamma\vee\alpha)=\beta$} / i.e. $C(\alpha)\wedge(\gamma\vee\alpha)=\gamma$, $C(\unsure{\beta|\gamma})(\gamma\vee\alpha)=\alpha$}
% slanted
\margin{\unsure{les} $\alpha,\beta,\gamma$ sont disjoints}
\end{enumerate}

\uline{Cor.} Conditions équivalentes sur $\Sigma$
\begin{enumerate}[label={},leftmargin=1.8em]
\item (i) $\forall$ couples $b,c$ avec $b\leq c$, les conditions du lemme 1 satisfaites.
\item (ii) $\forall$ couples $b,c$ avec $b\leq c$, les conditions \inserted{(plus fortes)} du lemme 2 sont satisfaites.
\end{enumerate}
% margin note written sideways
% slanted
\margin{\deleted{\unsure{ici}} \unsure{On peut voir|On peut remplacer} une famille finie $0\leq b_1\leq b_2\mathrel{\unsure{-|\cdots}}\leq b_n\leq1$ (pour faire (ii'))}
% margin note written sideways
% slanted
\margin{\unsure{dégager} \unsure{lemme} 3 \unsure{pour compléter}, \deleted{\unsure{articuler}} (\unsure{en} Cor. 2) \unsure{plus bas}}
\noindent$\left[\rule{0pt}{10em}\right.$\begin{minipage}{0.88\linewidth}
(iii) $\forall$ couples $\alpha,\beta$ avec $\alpha\,|\circ|\,\beta$, $\alpha$ et $\beta$ engendrent une sous-algèbre de Boole (laquelle sera \unsure{automatiquement} finie de cardinal $2^n$, $\deleted{1\leq n\leq}n\in\{1,2,3\}$), $n=1$ sss $\{\alpha,\beta\}\unsure{\in|\subset}\{0,1\}$, $n\leq2$ sss $\alpha$ ou $\beta$ est nul \inserted{ou $\alpha\vee\beta=1$} \deleted{\unsure{pour tous}} \deleted{sss \unsure{(NB)} $\alpha\ne0$, \unsure{et}} \deleted{\unsure{alors}} $n=3$ sss $\alpha\ne0$, $\beta\ne0$, $\alpha\vee\beta\ne1$.

\inserted{\unsure{(iiii)|(iiii')|(iii')|(iv)} Comme (iii), mais avec familles \uline{finies} $(\alpha_i)$, avec $\alpha_i\,|\circ|\,\alpha_j$ pour $i\ne j$}
\end{minipage}
